% ─────────────────────────────────────────────────────────────────────────────
% Wearable-Defined Exercise Episode Detection and Residual Audit
% Numbers sourced from:
%   notebooks/outputs/exercise_episode_detection_v2/detector_config.json
%   notebooks/outputs/exercise_episode_detection_v2/exercise_detection_summary.csv
%   notebooks/outputs/exercise_episode_detection_v2/onset_refinement_summary.csv
%   outputs/stage2_residual_diagnostic/diagnostic_verdict.json
%   outputs/stage2_residual_diagnostic/residual_comparison_summary.csv
%   outputs/stage2_residual_diagnostic/residual_trajectory_summary.csv
% No numbers invented.
% ─────────────────────────────────────────────────────────────────────────────
\section{ Exercise detection }
\label{sec:exercise_detection}

\subsection{Motivation}
\label{subsec:ex_motivation}

The meal counterfactual is not identifiable on AI-READI
(Section~\ref{sec:study1_cf_audit}), and the event aware residual correction
(Section~\ref{sec:study2}) a statistically significant but modest
improvement of $0.175$~mg/dL on triggered windows.
AI-READI participants wore Garmin wrist devices that record heart rate (HR),
step counts, and respiratory rate (RR) continuously, making ambulatory
exercise a second direction for counterfactuals.
Unlike meal events, physical activity produces large, time series physiological
signals in HR and steps that are detectable without any glucose information.

The goals: construct a detector of episodes that are 
ambulatory exercise using only wearable signals; characterize
the descriptive glucose response at detected onsets; and analyse the frozen
base model's residuals at episodes that are exrcices to determine whether systematic
signed bias justifies a correction or advisory.

No causal claims about exercise physiology are made, and all findings are
scoped to non-insulin participants unless noted.

\subsection{Available wearable wignals and design constraints}
\label{subsec:ex_signals}

The Garmin wrist stream provides HR (bpm), step counts (steps per 5-minute
epoch), RR (breaths/min), a sleep-stage label, and device-derived fields
including estimated calories per minute and a Garmin activity intensity score.

Four design constraints govern signal use in the detector.
First, glucose is never used to detect or gate episodes; glucose
appears only as an outcome variable.
Second, calories per minute and the Garmin activity intensity score
are excluded from all gates because both are partly derived from HR; using them
would circularly reinforce the HR criterion.
Third, a rolling local HR delta was considered as an onset refinement
feature but was reduced to visualization and onset diagnosis only because a
backward looking delta over exercise windows self-contaminates the anchor test.
Fourth, RR is used descriptively but not as a gate: its noise level
on wrist PPG is too high for a reliable onset criterion at 5-minute resolution.

The working signals for detection are therefore HR (excess above a stable
resting floor), step count rate, and the sleep label (for daytime restriction).


\label{subsec:ex_gate}

\paragraph{Stable resting HR floor:}
Each participant's resting HR floor is defined as the 5th percentile of
all HR observations recorded during sedentary, awake windows
(step rate below threshold, sleep label absent).
This is a person-specific baseline  computed once per
participant; it is never a rolling window, which would absorb exercise induced
HR elevation and suppress the excess signal.

\paragraph{HR excess: }
At each 5-minute anchor $t$, the HR excess above the stable floor is
\begin{equation}
  \Delta\mathrm{HR}_{t} = \mathrm{HR}_{t} - \mathrm{HR}_{\text{floor}},
  \label{eq:hr_excess}
\end{equation}
where $\mathrm{HR}_{\text{floor}}$ is the 5th-percentile stable resting HR
for that participant.

\paragraph{Anchor gate:}
A window is flagged as an anchor (potential exercise start) when all three
of the following hold simultaneously:
\begin{align}
  \overline{\mathrm{steps}}_{t,15\text{min}} &\geq 40 \text{ steps/min} \quad (\text{STEP\_RATE\_FLOOR}), \label{eq:gate_steps} \\
  \overline{\Delta\mathrm{HR}}_{t,15\text{min}} &\geq 8 \text{ bpm}    \quad (\text{HR\_EXCESS\_FLOOR}), \label{eq:gate_hr} \\
  \mathrm{sleep}_{t}                            &= 0.                  \label{eq:gate_wake}
\end{align}
Steps lead the gate; HR excess confirms physiological arousal.

\paragraph{Strain grading: }
Strain is graded against a person specific 20th-percentile HR excess floor
(not the detection gate 5th-percentile floor), creating three bands:
light ($\Delta\mathrm{HR} \in [8, 15)$~bpm), moderate ($[15, 30)$~bpm),
and vigorous ($\geq\!30$~bpm).
A sub threshold class captures anchors that clear the 5th percentile step
criterion but fall below the 8-bpm excess floor; these are retained in
cadence analysis only.


\label{subsec:ex_episodes}
\paragraph{Episode construction.}
Consecutive anchor windows are merged into sessions
(merge gap $\leq 5$~min, minimum episode duration $= 10$~min, maximum $= 120$~min).
Exercice episodes are excluded if they overlap with more than 5\% of
sleep stage annotations or form part of a long active envelope
($\geq 3$~hr of continuous activity), leaving the clean set (primary definition).
A stricter sensitivity subset additionally requires the post episode
HR-recovery fraction to exceed 30\% within 30 minutes
(strict recovery definition).

The clean set contains \textbf{2{,}171 episodes} across \textbf{788 participants}.
The strict recovery sensitivity subset contains \textbf{1{,}015 episodes}
across \textbf{574 participants}.

\paragraph{Confidence score and refined start of exercise:}
A confidence score weights five signals:
activity quality ($w = 0.40$), HR excess ($w = 0.30$),
RR elevation ($w = 0.10$, descriptive), duration quality ($w = 0.10$),
and recovery ($w = 0.10$); a sleep overlap penalty overrides all positive
weights ($w = -1.0$).
A high confidence episode exceeds a score of 0.65; medium confidence falls
between 0.40 and 0.65.
The confidence window start (when the episode first exceeds the anchor gate)
is not used as the onset: the detector backtracks from the
confidence start to the earliest sustained step rise within 45 minutes,
producing a \textbf{refined start} that precedes the confidence gate.

Across all 3{,}131 onset-refined bouts, the \textbf{median refined-onset
advance is 30 minutes} (mean 41.6~min) relative to the confidence-window
start.
Onset and offset are symmetric: if the start is pulled back, the
end is extended by the same logic, subject to a sleep-label clamp
and the 120-min duration cap.

\begin{figure}[H]
  \centering
  \includegraphics[width=0.95\linewidth]{exercise/aligned_ambulatory_confidence_vs_refined.pdf}
  \caption{%
    \textbf{Confidence-window onset versus refined onset.}
    Each episode is aligned to its confidence-window start (vertical dashed
    line at $t=0$); traces show mean step rate (top) and HR excess (bottom)
    for 90~min before and after that anchor.
    The median refined onset (30~min before the confidence start) captures
    the sustained step rise that precedes the confidence gate, confirming
    that the confidence-window start systematically lags the earlier
    wearable-defined activity signal.%
  }
  \label{fig:exercise_confidence_vs_refined}
\end{figure}

\subsection{Detector Validation}
\label{subsec:ex_validation}

\paragraph{Episode timelines: }
The first alignment below shows the per-episode
wearable signal alignment at the refined onset, pooled across all clean
episodes.
Steps rise sharply at $t=0$ and HR excess follows within 5--10 minutes,
consistent with an activity led gate in which step elevation precedes the
cardiovascular response.
The paired timelines below illustrate clean detection across an active day
for participants 1029 and 1030; participant 4022 is shown
as a diagnostic example with more complex episode structure and is discussed
in supplementary context only.

\begin{figure}[H]
  \centering
  \includegraphics[width=0.95\linewidth]{exercise/aligned_ambulatory_refined_start.pdf}
  \caption{%
    \textbf{Wearable signals aligned at the refined episode onset.}
    Mean and 95\% bootstrap CI for step rate (top) and HR excess
    above the stable resting floor (bottom), aligned at $t=0$
    (refined onset) across all 2{,}171 clean ambulatory exercise-like
    episodes.
    Steps rise sharply at onset; HR follows within 5--10 minutes,
    consistent with the activity-led gate design in which steps lead
    and HR confirms.
    No glucose signal is used in detection; the glucose trace visible
    in the full notebook is shown post-hoc only.%
  }
  \label{fig:exercise_refined_start}
\end{figure}

\begin{figure}[H]
  \centering
  \includegraphics[width=0.95\linewidth]{exercise/negative_control_validation.pdf}
  \caption{%
    \textbf{Negative control validation.}
    Three conditions aligned at a common reference: (i) real refined
    exercise like onsets; (ii) the same onsets time-shifted by
    $+6$~hours (temporal negative control); and (iii) random
    awake anchors matched on time-of-day (distributional negative control).
    Real onsets show a sharp, coordinated HR and step rise.
    Both negative controls are flat across the alignment window,
    validating that the detector captures a real physiological signal
    rather than circadian variation or device artefact.
    Exact mean HR-excess and step-rate values at $t=0$ for each condition
    are readable from the figure.%
  }
  \label{fig:exercise_negative_controls}
\end{figure}

Using the pre registered rise contrast, defined as the mean over
$[0,15]$ minutes minus the mean over $[-15,0]$ minutes, real onsets showed
an HR excess rise of $8.01$~bpm and a step-rate rise of
$29.44$~steps/min.  The $+6$~hour control showed $0.27$~bpm and
$-0.93$~steps/min, while the random-awake control showed $-0.46$~bpm and
$-0.23$~steps/min.  The resulting separation from the strongest control
was $7.74$~bpm for HR excess and $29.67$~steps/min for step rate.

\begin{figure}[H]
  \centering
  \begin{subfigure}[b]{0.49\linewidth}
    \centering
    \includegraphics[width=0.95\linewidth]{exercise/timeline_refined_1029_20230910.pdf}
    \caption{Participant 1029, 2023-09-10}
  \end{subfigure}
  \hfill
  \begin{subfigure}[b]{0.49\linewidth}
    \centering
    \includegraphics[width=0.95\linewidth]{exercise/timeline_refined_1030_20230910.pdf}
    \caption{Participant 1030, 2023-09-10}
  \end{subfigure}
  \caption{%
    \textbf{Representative episode timelines (participants 1029 and 1030).}
    Each panel shows a full day of HR, step rate, sleep label, and
    detected episode boundaries (refined onset/offset, shaded region).
    Both participants show clear HR-step coupling at episode boundaries.
    Participant 4022 (available as a supplementary figure) illustrates
    a participant with more complex episode structure and is used for
    diagnostic illustration only.%
  }
  \label{fig:exercise_timeline_examples}
\end{figure}

\subsection{Exercise type classification}
\label{subsec:ex_grading}

The episode type is characterized along two independent axes.
\textbf{Cadence} is the mean step rate over active core minutes
(minutes with $\geq 20$ steps/min within the episode) and is
person-invariant: it reflects movement pattern independent of
the participant's cardiovascular fitness.
Four cadence classes are defined: slow walk (40--80 steps/min),
walk (80--110), brisk walk (110--130), and run-like ($\geq\!130$).
``Run-like'' denotes high-cadence ambulation and does not confirm
a specific activity mode.

\textbf{Strain} is the HR excess above the person-specific 20th-percentile
sedentary-awake floor and is person-specific: it reflects cardiovascular
load relative to that participant's own baseline.
Three strain bands are defined: light (8--15~bpm excess), moderate
(15--30~bpm), and vigorous ($\geq\!30$~bpm).

\begin{figure}[H]
  \centering
  \includegraphics[width=0.95\linewidth]{exercise/cadence_strain_heatmap.pdf}
  \caption{%
    \textbf{Cadence--strain episode heatmap.}
    Each cell shows episode count (or a derived metric) for the
    cadence-class (columns) by strain class (rows) joint distribution
    across the 2{,}171 clean episodes.
    Slow walk and walk cadence classes dominate numerically.
    A notable feature is that many slow-walk episodes carry moderate or
    vigorous strain, reflecting inter-individual variability: a slow gait
    can be cardiovascularly demanding for participants with low fitness or
    high resting HR variability.
    This heterogeneity motivates stratified analysis later
    (Section~\ref{subsec:ex_residual}) but does not support inference about
    insulin resistance or cardiovascular fitness.%
  }
  \label{fig:exercise_cadence_strain}
\end{figure}

\subsection{Recovery and Quality Control}
\label{subsec:ex_recovery}

HR recovery is assessed as the fraction of the HR excess that dissipates
within 30 minutes post episode (recovery\_frac\_thresh $= 0.30$).
Recovery is used as a \emph{sensitivity marker}, not a hard detection
criterion: episodes failing recovery are retained in the primary set but
excluded from the strict-recovery sensitivity subset.
This design avoids circularity in which a recovery criterion would
preferentially select short, high-intensity bouts that happen to have
large post-episode HR drops.

\begin{figure}[H]
  \centering
  \includegraphics[width=0.95\linewidth]{exercise/recovery_aligned_end.pdf}
  \caption{%
    \textbf{HR recovery aligned at episode end.}
    Mean HR excess (with 95\% CI) in the 60~minutes following the
    episode end, stratified by strict-recovery pass/fail.
    Episodes that pass the strict-recovery criterion (fraction $\geq 0.30$
    within 30~min post-end) show a clear HR return toward baseline;
    episodes that fail show a slower or partial return, consistent with
    extended moderate-activity transitions rather than detector error.
    Recovery rate is interpreted as a sensitivity marker for episode
    quality, not a physiological claim.%
  }
  \label{fig:exercise_recovery}
\end{figure}

\subsection{Preliminary glucose-response kernel}
\label{subsec:ex_glucose_kernel}

The aligned kernel below shows mean glucose at the refined episode onset,
pooled across all clean episodes with available CGM.
The descriptive profile shows a \textbf{16.66~mg/dL drop}
(the moderate-strain, non-insulin anchor used in the planning response
in Section~\ref{subsec:ex_pathb}), with a minimum reached at
\textbf{28.31 minutes}
post onset, followed by a partial rebound within the observation window.

This is a \textbf{purely descriptive, associational pattern}.
It is not a causal estimate.
The observed trajectory is potentially confounded by meals occurring during or
near exercise starts, pre exercise glucose level (regression to the mean),
circadian time-of-day, and selection of episode type.
Insulin treated participants are excluded from causal interpretation throughout;
the kernel here is descriptive across all participants with CGM coverage.
The kernel is later used as a magnitude anchor for the planning response
(Section~\ref{subsec:ex_pathb}), where this descriptive-associational
role is stated explicitly.

\begin{figure}[H]
  \centering
  \includegraphics[width=0.95\linewidth]{exercise/glucose_aligned_at_onset.png}
  \caption{%
    \textbf{Descriptive glucose kernel aligned at refined episode onset.}
    Mean glucose (mg/dL) with 95\% bootstrap CI, aligned at $t=0$
    (refined onset) across clean episodes with CGM coverage.
    The profile shows an approximate 16.7~mg/dL drop and a partial rebound.
    This is a descriptive, associational pattern: confounding by meals,
    baseline glucose, and time-of-day is possible and not controlled.
    Glucose is never used to detect episodes; this is a post-hoc summary.%
  }
  \label{fig:exercise_glucose_kernel}
\end{figure}

\subsection{Comparison to the frozen base model }
\label{subsec:ex_residual}

The parts focuses on whether the frozen \texttt{mamba\_stateful\_5epoch}
base model (the same checkpoint as Study~2) makes systematically larger errors
inside exercise like windows than outside.
For each clean episode with forecast cache coverage, a matched non-exercise
control window is drawn from the same participant, matched within
$\pm 60$~min time-of-day and $\pm 20$~mg/dL baseline glucose.
The signed residual follows the notebook convention:
\begin{equation}
  b_{h} = \hat{q}_{50,t}^{(h)} - G_{t+h},
  \label{eq:signed_residual}
\end{equation}
where $\hat{q}_{50,t}^{(h)}$ is the frozen base median forecast at horizon $h$
issued from anchor $t$.
Positive $b_{h}$ indicates that the base model over predicts actual glucose,
so the observed glucose is lower than forecast.  The horizon-growing positive
bias therefore indicates that the frozen model under represents the
exercise associated drop.

Matching yielded \textbf{212 episode--control pairs}
($316$ episodes passed Step~0, $271$ had cache coverage, and $212$ found a
within-participant matched control at the selected tolerance).

\paragraph{MAE comparison.}
Episode and control windows show no statistically distinguishable MAE gap:
paired MAE gap (all strains): $+0.29$~mg/dL, 95\% CI $[-0.68,\;+1.28]$
(CI includes zero).
The null MAE finding holds within each strain class
(vigorous: $+0.67$ $[-0.51,\;+1.83]$;
moderate: $-0.17$ $[-1.72,\;+1.56]$).
The frozen model is not less accurate during exercise like windows on the
absolute error metric.

\paragraph{Signed bias }
The signed bias reveals a different picture.
The model's forecast progressively over predicts the actual glucose
trajectory during exercise-episode horizons compared to matched controls,
growing with forecast horizon.
At the matched-pairs level:

\begin{itemize}
  \item \textbf{Vigorous} ($n = 104$ pairs):
        episode signed bias $+3.82$ mg/dL (95\% CI $[+2.35,\;+5.21]$);
        control signed bias $-0.37$ mg/dL ($[-1.51,\;+0.68]$);
        \textbf{paired bias gap $+4.19$ mg/dL} (CI $[+2.33,\;+5.97]$,
        CI excludes zero).
  \item \textbf{Moderate} ($n = 102$ pairs):
        episode signed bias $+2.20$ mg/dL (95\% CI $[+0.09,\;+4.09]$);
        control signed bias $-1.85$ mg/dL ($[-3.40,\;-0.33]$);
        \textbf{paired bias gap $+4.05$ mg/dL} (CI $[+1.83,\;+6.29]$,
        CI excludes zero).
  \item \textbf{Light} ($n = 6$ pairs):
        paired bias gap $+4.43$ mg/dL (CI $[-4.47,\;+15.11]$,
        CI includes zero). \emph{Do not interpret: very low sample size.}
\end{itemize}

\noindent
The bias gap for vigorous and moderate strain both exclude zero, indicating
provisional, recoverable headroom above matched controls.
The verdict from the diagnostic pipeline is
\textbf{JUSTIFIED} (vigorous headroom exceeds the 2~mg/dL material-gap
threshold; moderate also exceeds it but with a wider CI).
The headroom is termed \emph{provisional} here because:
the analysis is scoped to non-insulin T2D participants (descriptive for
insulin-treated; the insulin-treated subgroup has $n = 32$ pairs and
wide CIs); episode count is modest; matching tolerance is coarse.

\begin{figure}[H]
  \centering
  \includegraphics[width=0.95\linewidth]{exercise/episode_vs_control_residual_by_strain.png}
  \caption{%
    \textbf{Episode versus matched-control signed residual by strain class.}
    Each panel shows the mean signed residual (mg/dL) for exercise-like
    episodes (red) versus matched non-exercise controls (blue) at the
    median horizon, stratified by strain class.
    Error bars are 95\% bootstrap CIs.
    Vigorous and moderate episodes show a positive bias gap relative to
    matched controls (CIs excluding zero), while light episodes
    ($n = 6$ pairs) show a wide CI.
    No causal claim is made: the bias reflects the frozen model's inability
    to anticipate the exercise-associated glucose recovery, not a confirmed
    exercise effect.%
  }
  \label{fig:exercise_residual_strain}
\end{figure}

\begin{figure}[H]
  \centering
  \includegraphics[width=0.95\linewidth]{exercise/residual_trajectory_by_horizon.png}
  \caption{%
    \textbf{Signed residual trajectory by forecast horizon and strain class.}
    Mean signed residual ($b_h = \hat{q}_{50}^{(h)} - G_{t+h}$, mg/dL) for
    episodes (red) and matched controls (blue) as a function of forecast
    horizon (5--60 min), stratified by strain class (light, moderate, vigorous).
    The episode residual grows monotonically with horizon for moderate and
    vigorous strain, while the control residual remains near zero.
    Light episodes ($n = 6$) show a wide CI that precludes interpretation.
    This horizon-growing bias is the provisional headroom that motivates
    both the direction of planning response(counterfactual) and advisory correction  in
    Section~\ref{sec:exercise_counterfactual}.%
  }
  \label{fig:exercise_residual_horizon}
\end{figure}

\subsection{Limitations bridging to the counterfactual}
\label{subsec:ex_limitations}

The detector is wearable defined, not ground truth validated against
prospective exercise logs.
Detection captures only ambulatory, step based activity; low step activities such as resistance training, isometric exercise, stationary cycling are not 
captured.
RR is too noisy for reliable onset detection at wrist-PPG resolution and is
used descriptively only.
Meal confounding is possible within the glucose kernel: episode windows
of 10--120 min can overlap with eating, particularly at mealtimes.

The residual analysis is scoped to non insulin participants for causal
interpretation; insulin treated results are reported descriptively
because unobserved insulin dosing is a direct confounder of the glucose
trajectory.
Subgroup analyses by cadence class, time-of-day, and participant phenotype
are deferred until the detector and kernel are confirmed stable on a
holdout population.

The insulin resistance hypothesis was tested and closed as a negative result:
baseline glucose was the only predictor of the subsequent drop, and no
recoverable participant exercise sensitivity term was found.  The relationship
between cadence/strain heterogeneity and participant metabolic profiles is
therefore framed as a fairness and heterogeneity issue to characterize in
future work, not as evidence for a biological mechanism.

% ── Summary Table ─────────────────────────────────────────────────────────────

\begin{table}[H]
  \centering
  \caption{%
    \textbf{EX-T1: Detector and residual audit summary.}
    All values are sourced from notebook or pipeline output artifacts.%
  }
  \label{tab:exercise_summary}
  \small
  \setlength{\tabcolsep}{3pt}
  \begin{tabular}{p{0.20\linewidth}p{0.37\linewidth}p{0.34\linewidth}}
    \toprule
    Item & Result & Interpretation \\
    \midrule
    Clean episodes (primary)    & 2{,}171           & Primary detection set \\
    Participants (clean)        & 788               & Unique participants with $\geq 1$ clean episode \\
    Strict-recovery subset      & 1{,}015 episodes  & 574 participants; sensitivity check \\
    Median onset shift          & 30 min            & Refined onset precedes confidence-window start \\
    Negative-control HR rise    & 8.01 vs 0.27 / $-0.46$ bpm & Real vs shifted / random-awake \\
    Negative-control steps rise & 29.44 vs $-0.93$ / $-0.23$ steps/min & Real vs shifted / random-awake \\
    Cadence distribution        & 1{,}859 slow; 223 walk; 81 sub-threshold; 6 brisk; 2 run-like & Slow walk dominates \\
    Glucose kernel drop         & 16.66 mg/dL & Descriptive-associational; anchor for Path~B \\
    Glucose kernel nadir        & 28.31 min & Moderate-strain, non-insulin anchor \\
    MAE gap (vigorous)          & $+0.67$ $[-0.51,\;+1.83]$~mg/dL & CI includes zero; no error increase \\
    MAE gap (moderate)          & $-0.17$ $[-1.72,\;+1.56]$~mg/dL & CI includes zero; no error increase \\
    Signed bias gap (vigorous)  & $+4.19$ $[+2.33,\;+5.97]$~mg/dL & CI excludes zero; provisional headroom \\
    Signed bias gap (moderate)  & $+4.05$ $[+1.83,\;+6.29]$~mg/dL & CI excludes zero; provisional headroom \\
    Residual audit verdict      & JUSTIFIED         & Vigorous headroom $>$ 2 mg/dL threshold \\
    \bottomrule
  \end{tabular}
\end{table}
